%! AP Statistics — Unit 9, Lesson 9.3 — Lesson Plan
\documentclass[10pt]{article}
\usepackage{apstats-article}
\usepackage{apstats-boxes}
\usepackage{graphicx}
\graphicspath{{images/}}
\newcommand{\UnitNumberName}{Unit 9: Inference for Quantitative Data: Slopes \quad}
\newcommand{\LessonNumberName}{Lesson 9.3: Justifying a Claim About the Slope of a Regression Model Based on a Confidence Interval}

\begin{document}
\begin{center}
    {\Large\bfseries \CourseName: \SchoolYear \\
    {\small\bfseries \UnitNumberName \LessonNumberName}}
\end{center}

\begin{tcolorbox}[colback=sky, colframe=navy, boxrule=0.9pt, arc=2mm,
    left=2mm, right=2mm, top=1.5mm, bottom=1.5mm]
    \textbf{Primary Objective:} Students will interpret a confidence interval for the slope
    of a regression model in context, use a CI to justify or refute a claim about the true slope
    $\beta$, and describe how sample size affects interval width. (Big Idea: UNC)
\end{tcolorbox}

\renewcommand{\arraystretch}{1.6}
\begin{skillbox}[Priority Ideas \& Skills]{goldbox}
\begin{minipage}[t]{0.48\linewidth}
    \textbf{AP Skills Addressed}
    \begin{itemize}
        \item Interpret a CI for the slope of a regression model, referencing the sample and
              population with correct units (Skill 4.B; UNC-4.AG).
        \item Justify a claim based on whether a specific value falls inside or outside a CI
              for the slope (Skill 4.D; UNC-4.AH).
        \item Identify the effect of increasing sample size on the width of a CI for the slope
              (Skill 4.A; UNC-4.AI).
    \end{itemize}
\end{minipage}
\hfill
\begin{minipage}[t]{0.48\linewidth}
    \textbf{Key Understandings}
    \begin{itemize}
        \item Interpretation template: ``We are C\% confident that the true slope of the
              population regression line relating [y] to [x] for [population] is between
              [lower] and [upper] [units].''
        \item If 0 is NOT in the CI $\Rightarrow$ evidence that $\beta \ne 0$ (real linear
              relationship).
        \item If 0 IS in the CI $\Rightarrow$ consistent with $\beta = 0$; cannot conclude
              a linear relationship exists.
        \item Larger $n$ $\Rightarrow$ smaller $SE_b$ $\Rightarrow$ narrower CI, all else equal
              (UNC-4.AI.1).
    \end{itemize}
\end{minipage}
\end{skillbox}

\begin{skillbox}[{Vocabulary, Concepts, \& Theorems}]{greenbox}
\renewcommand{\arraystretch}{1.3}
\begin{tabularx}{\linewidth}{lX}
    \textbf{CI for slope} &
        An interval estimate of the true slope $\beta$; form $b \pm t^* \cdot SE_b$
        (UNC-4.AG.1--2) \\
    \textbf{Interpretation of CI for slope} &
        Must include: C\% confidence level, both endpoints with units, and a reference to the
        \emph{population} (not just the sample) (UNC-4.AG.2) \\
    \textbf{Justifying a claim with a CI} &
        If the claimed value is outside the CI $\Rightarrow$ evidence against that claim;
        if inside $\Rightarrow$ consistent with the data (UNC-4.AH.1) \\
    \textbf{Width of CI for slope} &
        Width $= 2t^* \cdot SE_b$; larger $n$ decreases $SE_b$, producing a narrower interval
        when all other things remain the same (UNC-4.AI.1) \\
\end{tabularx}
\end{skillbox}

\begin{skillbox}[Activate Prior Knowledge \& Spiral Review]{sky}
\begin{minipage}[t]{0.48\linewidth}
    \textbf{Prior Knowledge Check (3--4 min)}
    \begin{itemize}
        \item Lesson 9.2: constructing a CI for slope; the roles of $b$, $SE_b$, and $t^*$.
        \item Unit 7: interpreting a CI for a mean --- ``We are C\% confident that the true
              mean\ldots''
        \item Lesson 9.1: the LINE-R conditions for inference on the slope.
    \end{itemize}
\end{minipage}
\hfill
\begin{minipage}[t]{0.48\linewidth}
    \textbf{Connections Forward}
    \begin{itemize}
        \item Significance testing for slope $\to$ Lesson 9.4
        \item Connecting CI conclusions to $p$-value conclusions $\to$ Lesson 9.5
        \item AP free-response writing for regression inference $\to$ Unit 9 review
    \end{itemize}
\end{minipage}
\end{skillbox}

\begin{skillbox}[Hook \& Lesson]{sky}
\begin{minipage}[t]{0.48\linewidth}
    \textbf{Hook (4--5 min)}\\
    Display computer output: a 95\% CI for the slope relating daily study time (hours) to
    exam score (points) for high school students is $(0.14,\ 0.52)$.\\[4pt]
    Ask: ``A researcher claims the true slope is 0 --- students who study more don't actually
    score higher. Can you use this interval to evaluate that claim? What if the claim was
    $\beta = 0.5$?''\\[4pt]
    Let students discuss briefly, then frame: today's lesson is about reading what a CI for
    slope \emph{says} about a population and using it to evaluate specific claims.
\end{minipage}
\hfill
\begin{minipage}[t]{0.48\linewidth}
    \textbf{Lesson Flow (15 min)}
    \begin{enumerate}
        \item Model the interpretation template; annotate each required element on the board.
        \item Introduce claim-justification logic: locate the claimed value relative to the
              interval; outside = evidence against the claim.
        \item Special case: 0 not in CI $\Rightarrow$ evidence that $\beta \ne 0$; 0 in CI
              $\Rightarrow$ cannot conclude a linear relationship.
        \item Width and sample size: show that $SE_b$ decreases as $n$ increases, so larger
              samples produce narrower CIs, all else equal.
        \item AP-writing annotation: compare a full-credit model response to a common
              incomplete response.
    \end{enumerate}
\end{minipage}
\end{skillbox}

\begin{skillbox}[Explicit Instruction: Interpreting and Using a CI for Slope]{sky}
\begin{minipage}[t]{0.48\linewidth}
    \textbf{Interpretation Template}\\[4pt]
    ``We are \underline{C\%} confident that the true slope of the population regression line
    relating \underline{[response]} to \underline{[explanatory]} for \underline{[population]}
    is between \underline{[lower]} and \underline{[upper]} \underline{[units]}.''\\[6pt]
    \textbf{Required elements for full AP credit:}
    \begin{enumerate}
        \item Confidence level (C\%)
        \item Both endpoints with correct units
        \item Reference to the \emph{population} (not the sample)
    \end{enumerate}
\end{minipage}
\hfill
\begin{minipage}[t]{0.48\linewidth}
    \textbf{Claim-Justification Logic}\\[4pt]
    Given 95\% CI for $\beta$: $(L,\ U)$
    \begin{itemize}
        \item \textbf{Claim: $\beta = 0$}\\
              $0 \notin (L,U)$ $\Rightarrow$ evidence that $\beta \ne 0$;
              the data support a linear relationship.\\
              $0 \in (L,U)$ $\Rightarrow$ consistent with no linear relationship.
        \item \textbf{Claim: $\beta = \beta_0$}\\
              $\beta_0 \notin (L,U)$ $\Rightarrow$ evidence against that claim.\\
              $\beta_0 \in (L,U)$ $\Rightarrow$ claim is plausible.
        \item \textbf{Width}: $n \uparrow\ \Rightarrow SE_b \downarrow\ \Rightarrow$
              narrower CI.
    \end{itemize}
\end{minipage}
\end{skillbox}

\begin{skillbox}[Active Monitoring]{sky}
\begin{minipage}[t]{0.48\linewidth}
    \textbf{Circulate and Check}
    \begin{itemize}
        \item Students often omit the \emph{population} reference --- redirect to ``for all
              [population],'' not just the sample.
        \item Watch for students saying ``there IS a relationship'' when 0 is in the CI;
              correct to ``the data are consistent with no relationship.''
        \item Ensure units of slope are reported correctly (response units per explanatory unit).
    \end{itemize}
\end{minipage}
\hfill
\begin{minipage}[t]{0.48\linewidth}
    \textbf{Cold-Call Prompts}
    \begin{itemize}
        \item ``Why must the interpretation reference the \emph{population}, not just the
              sample?''
        \item ``A 95\% CI for slope is $(-0.3,\ 0.8)$. Does this provide evidence that
              $\beta \ne 0$? Explain.''
        \item ``If a researcher doubles the sample size and recomputes the CI, what happens
              to the width? Why?''
    \end{itemize}
\end{minipage}
\end{skillbox}

\begin{skillbox}[Group Work \& Differentiation]{redbox}
\begin{minipage}[t]{0.48\linewidth}
    \textbf{Partner Activity (10--12 min)}\\
    Students receive three regression computer outputs (hours studied vs.\ exam score;
    outside temperature vs.\ daily ice cream sales; age vs.\ resting heart rate) with 95\%
    CIs already computed. For each: interpret the CI, evaluate a researcher's specific claim,
    and compare interval widths across different sample sizes. Three tiers (R, A, E).
\end{minipage}
\hfill
\begin{minipage}[t]{0.48\linewidth}
    \textbf{Differentiation}
    \begin{itemize}
        \item \textit{Support (Tier R)}: Sentence frame provided; circle ``is'' or ``is not''
              for claim evaluation.
        \item \textit{On-grade (Tier A)}: Write full interpretations, justify claims, compare
              widths.
        \item \textit{Extension (Tier E)}: Also explain why width changes via $SE_b$; critique
              a flawed student interpretation.
        \item \textit{ELL}: Vocabulary support --- ``slope / pendiente,'' ``confident /
              confiado,'' ``claim / afirmaci\'{o}n.''
    \end{itemize}
\end{minipage}
\end{skillbox}

\begin{skillbox}[Individual Work \& Assessment]{redbox}
\begin{minipage}[t]{0.48\linewidth}
    \textbf{Exit Ticket (5 min)}
    \begin{enumerate}
        \item Given a 95\% CI for slope in context, write a complete interpretation.
        \item Does the interval provide evidence that $\beta \ne 0$? Justify.
        \item A researcher claims $\beta = 0.5$ --- is this consistent with the interval?
    \end{enumerate}
\end{minipage}
\hfill
\begin{minipage}[t]{0.48\linewidth}
    \textbf{AP-Style Check}\\
    \textit{A 95\% CI for the slope relating weight (kg) to height (cm) for adult males is
    $(0.48,\ 0.92)$. Which statement is best supported?}
    \begin{enumerate}[label=(\Alph*)]
        \item There is no evidence of a linear relationship.
        \item The true slope is exactly 0.70 kg/cm.
        \item The data provide evidence that the true slope is positive.
        \item We are 95\% confident that the \emph{sample} slope is between 0.48 and 0.92.
    \end{enumerate}
    Answer: (C)
\end{minipage}
\end{skillbox}

\begin{skillbox}[Reinforcement \& Extension]{goldbox}
\begin{minipage}[t]{0.48\linewidth}
    \textbf{Homework / Practice}
    \begin{itemize}
        \item Problem 1: given a computer output with a 95\% CI for slope, write a complete
              interpretation and evaluate two researcher claims (one where $\beta_0$ is inside
              the CI, one where it is outside).
        \item Problem 2: compare two CIs for slope from the same context with different sample
              sizes ($n = 20$ vs.\ $n = 80$); explain which is narrower and why.
    \end{itemize}
\end{minipage}
\hfill
\begin{minipage}[t]{0.48\linewidth}
    \textbf{Extension / Preview}
    \begin{itemize}
        \item \textit{Extension}: If a CI for slope contains only positive values, what does
              that imply about the $p$-value for a two-sided significance test at the same
              $\alpha$ level?
        \item \textit{Preview}: Lesson 9.4 introduces the $t$-test for slope --- students
              will see that the CI and the hypothesis test lead to the same conclusion.
    \end{itemize}
\end{minipage}
\end{skillbox}

\end{document}
